% ============================================================
%  Template Beamer - Algoritmos Computacionais I (UERJ)
%  Uso: copie este arquivo, altere \titulo / \subtitulo / \data
%  e escreva os slides entre \begin{document} e \end{document}.
%  Compilar com: pdflatex (ou latexmk -pdf)
% ============================================================
\documentclass[aspectratio=169,11pt]{beamer}

\usepackage[utf8]{inputenc}
\usepackage[T1]{fontenc}
\usepackage[brazil]{babel}
\usepackage{graphicx}
\usepackage{xcolor}
\usepackage{listings}

% ---------- Dados da aula (EDITE AQUI) ----------
\newcommand{\autor}{João Araujo}
\newcommand{\disciplina}{Algoritmos Computacionais I}
\newcommand{\titulo}{Título da Aula}
\newcommand{\subtitulo}{Subtítulo opcional}
\newcommand{\data}{\today}
\newcommand{\instituicao}{Departamento de Engenharia de Sistemas e Computação}

% Caminho do logo (relativo a este arquivo; ajuste se o .tex ficar em subpasta)
\newcommand{\logouerj}{imagens/logo_uerj_cor_transparencia.png}
% O logo é opcional: sem o arquivo ele é omitido.
\newcommand{\logoinst}[1]{\IfFileExists{\logouerj}{\includegraphics[#1]{\logouerj}}{}}
% Logo da capa: a quebra de linha só existe se o logo existir.
\newcommand{\logocapa}[1]{\IfFileExists{\logouerj}{\includegraphics[height=#1]{\logouerj}\\[1.2ex]}{}}

% ---------- Cores ----------
\definecolor{uerjazul}{RGB}{0,62,126}
\definecolor{uerjcinza}{RGB}{240,242,246}

% ---------- Tema ----------
\usetheme{default}
\usecolortheme{default}
\usefonttheme{professionalfonts}
\setbeamertemplate{navigation symbols}{}

\setbeamercolor{structure}{fg=uerjazul}
\setbeamercolor{frametitle}{fg=uerjazul,bg=white}
\setbeamercolor{title}{fg=uerjazul}
\setbeamercolor{block title}{fg=white,bg=uerjazul}
\setbeamercolor{block body}{bg=uerjcinza}
\setbeamercolor{footline}{fg=white,bg=uerjazul}
\setbeamercolor{itemize item}{fg=uerjazul}
\setbeamercolor{enumerate item}{fg=uerjazul}

\setbeamerfont{frametitle}{series=\bfseries,size=\large}
\setbeamerfont{footline}{size=\scriptsize}

% Título do frame com logo à direita
\setbeamertemplate{frametitle}{%
  \nointerlineskip
  \begin{beamercolorbox}[wd=\textwidth,ht=4ex,dp=1.2ex]{frametitle}
    \vspace{0.6ex}
    \usebeamerfont{frametitle}\insertframetitle%
    \hfill\raisebox{-1.2ex}{\logoinst{height=5ex}}
  \end{beamercolorbox}%
  \nointerlineskip
  {\color{uerjazul}\rule{\textwidth}{1.2pt}}%
}

% Rodapé: nome do autor | disciplina | número da página
\setbeamertemplate{footline}{%
  \leavevmode%
  \hbox{%
    \begin{beamercolorbox}[wd=.4\paperwidth,ht=2.5ex,dp=1.2ex,left,leftskip=1em]{footline}
      \usebeamerfont{footline}\autor
    \end{beamercolorbox}%
    \begin{beamercolorbox}[wd=.4\paperwidth,ht=2.5ex,dp=1.2ex,center]{footline}
      \usebeamerfont{footline}\disciplina
    \end{beamercolorbox}%
    \begin{beamercolorbox}[wd=.2\paperwidth,ht=2.5ex,dp=1.2ex,right,rightskip=1em]{footline}
      \usebeamerfont{footline}\insertframenumber{} / \inserttotalframenumber
    \end{beamercolorbox}%
  }%
}

% Página de título personalizada
\setbeamertemplate{title page}{%
  \vfill
  \begin{center}
    \logocapa{3.2cm}
    {\scriptsize\color{uerjazul}\instituicao}\\[2.5ex]
    {\usebeamerfont{title}\usebeamercolor[fg]{title}\Large\bfseries\inserttitle\par}
    \vspace{0.6ex}{\normalsize\color{gray}\insertsubtitle\par}
    \vspace{2ex}
    {\color{uerjazul}\rule{0.5\textwidth}{0.6pt}}\\[1.5ex]
    {\normalsize\autor}\\[0.6ex]
    {\small\disciplina}\\[0.6ex]
    {\footnotesize\data}
  \end{center}
  \vfill
}

% Metadados
\title{\titulo}
\subtitle{\subtitulo}
\author{\autor}
\institute{\instituicao}
\date{\data}

% Estilo para código
\lstset{
  basicstyle=\ttfamily\small,
  backgroundcolor=\color{uerjcinza},
  keywordstyle=\color{uerjazul}\bfseries,
  commentstyle=\color{gray}\itshape,
  stringstyle=\color{teal},
  frame=single,
  rulecolor=\color{uerjazul},
  showstringspaces=false,
  inputencoding=utf8,
  extendedchars=true,
  literate={á}{{\'a}}1 {é}{{\'e}}1 {í}{{\'i}}1 {ó}{{\'o}}1 {ú}{{\'u}}1
           {ã}{{\~a}}1 {õ}{{\~o}}1 {ç}{{\c c}}1 {â}{{\^a}}1 {ê}{{\^e}}1 {Ã}{{\~A}}1
}

% ============================================================
\begin{document}

% Capa: só título da aula (sem rodapé)
{
\setbeamertemplate{footline}{}
\begin{frame}[plain]
  \titlepage
\end{frame}
}

\begin{frame}{Roteiro}
  \tableofcontents
\end{frame}

\section{Introdução}
\begin{frame}{Exemplo de slide}
  \begin{itemize}
    \item Primeiro tópico
    \item Segundo tópico
  \end{itemize}

  \begin{block}{Definição}
    Texto em destaque.
  \end{block}
\end{frame}

\begin{frame}[fragile]{Exemplo de código}
\begin{lstlisting}[language=Python]
def soma(a, b):
    # soma dois números
    return a + b

print(soma(2, 3))
\end{lstlisting}
\end{frame}

\end{document}
